\section{Vacuum doubling with a weak asymptotic tail}
\label{uniq:section}

The vacuum pair obtained in Proposition~\ref{stat:vacuum-reduction}
has only the asymptotic derivative bounds inherited from the original
data. We prove its uniqueness at that regularity. In particular, the
point added in conformal doubling will initially be a
$W^{2,s}$ point, rather than a smooth point. The nonnegative mass
statement used below follows from the numerical
Theorem~\ref{thm:numerical}, which has already been proved; no equality
conclusion of that theorem is used.

\begin{proposition}[Vacuum uniqueness]\label{uniq:vacuum-uniqueness}
Let $n\geq4$ and let $(B,h_0,\lambda_0)$ be a smooth oriented vacuum static
exterior, complete including its connected compact boundary $S$,
with a compact complement of one Euclidean coordinate end. Suppose
\begin{equation}\label{uniq:static-equations}
 \Hess_{h_0}\lambda_0=\lambda_0\Ric_{h_0},\qquad
 \Delta_{h_0}\lambda_0=0,\qquad
 \lambda_0|_S=0,\qquad \lambda_0\longrightarrow1,
\end{equation}
and $\partial_\nu\lambda_0=\kappa_0>0$ on $S$, with $\nu$ pointing
toward the end. Assume, after a constant linear normalization of end
coordinates,
\begin{equation}\label{uniq:initial-tail}
 h_0-\delta=O_2(r^{-q_0}),\qquad
 \lambda_0-1=O_2(r^{-q_0}),\qquad
 \frac{n-2}{2}<q_0<n-2.
\end{equation}
Then $B$ is globally the exterior of a round sphere in isotropic
coordinates. For some $R>0$, writing $\rho=|x|\geq R$,
\begin{equation}\label{uniq:schwarzschild-pair}
 h_0=\left(1+\frac{R^{n-2}}{\rho^{n-2}}\right)^{4/(n-2)}\delta,
 \qquad
 \lambda_0=
 \frac{1-R^{n-2}/\rho^{n-2}}{1+R^{n-2}/\rho^{n-2}}.
\end{equation}
The identification extends smoothly to $S$.
\end{proposition}

We first give the end and analytic ingredients needed in its proof.

\begin{lemma}[Harmonic coordinates and the missing metric monopole]
\label{uniq:end-expansion}
Under the hypotheses of Proposition~\ref{uniq:vacuum-uniqueness}, put
\[
 \check h=\lambda_0^{2/(n-2)}h_0,
 \qquad V_0=\log\lambda_0.
\]
There are harmonic end coordinates for $\check h$, a number $a_0>0$,
and $\alpha>0$ such that, for every derivative order $j$,
\begin{equation}\label{uniq:improved-tail}
 \check h_{ij}=\delta_{ij}+O_j(r^{2-n-\alpha}),\qquad
 V_0=-a_0r^{2-n}+O_j(r^{2-n-\alpha}).
\end{equation}
\end{lemma}

\begin{proof}
The maximum principle gives $0<\lambda_0<1$ in the interior. The
conformal Ricci and Laplace formulas applied to
\eqref{uniq:static-equations} give
\begin{equation}\label{uniq:harmonic-ricci}
 \Ric_{\check h}=\frac{n-1}{n-2}\,\dd V_0\otimes\dd V_0,
 \qquad \Delta_{\check h}V_0=0.
\end{equation}
We may decrease $q_0$ while preserving its strict lower bound. Since
$n\geq4$, it is possible to take $q_0>1$; hence
$\delta_0=1-q_0$ lies in $(2-n,0)$. This avoids a logarithmic
coordinate correction.

Here are details of constructing coordinates at the stated derivative
count. Extend the coefficients of $\check h$ to be Euclidean inside
a large coordinate sphere, using a transition annulus. The equation
for $y^i=x^i+v^i$ is
\[
 \Delta_\delta v^i=-\Delta_{\check h}x^i
                 +(\Delta_\delta-\Delta_{\check h})v^i.
\]
Use weighted scaled $C^{2,\beta}$ norms of order $\delta_0$,
with a fixed $0<\beta<1$. The Newton operator takes order
$\delta_0-2$ to order $\delta_0$: splitting its integral into
$|z|<r/2$, $|z-x|<r/2$, and the remaining exterior annuli bounds the
undifferentiated potential by $C(1+r)^{\delta_0}$, and local scaled
Schauder estimates give its two derivatives and H\"older seminorm.
Both integrals at infinity converge because $\delta_0<0$, and
$\delta_0>2-n$ controls the inner annuli. After choosing the cutoff
radius sufficiently large, the last operator has norm less than one.
The Neumann series therefore gives
$v^i=O_{2,\beta}(r^{1-q_0})$. The original $O_2$ bounds supply the
scaled H\"older bounds for the first-order coefficients used here.
The new coordinate derivative is $I+O(r^{-q_0})$, so these are
coordinates outside a larger sphere. There is a derivative-count
issue in transferring a metric to these coordinates: its second
derivatives use third derivatives of the coordinate correction.
Differentiating the coordinate equation once, the assumed $O_2$
metric bounds control its right side in scaled $L^t$ for every
finite $t$. Local estimates give $v^i$ in scaled $W^{3,t}$ with
the corresponding decay. Thus the transformed metric first has
the initial weighted $W^{2,t}$ bounds for any finite $t$, which
is enough for the next elliptic step. No third asymptotic
derivative of the original metric has been assumed.

In these coordinates \eqref{uniq:harmonic-ricci} is a uniformly
elliptic system whose metric equation has principal part
$-\check h^{ab}\partial_a\partial_b\check h_{ij}/2$ and quadratic
first-derivative terms. The equation for $V_0$ has the same principal
coefficients. Scaled interior estimates first give $C^{2,\beta}$
bounds and then, by differentiating the system, all higher scaled
bounds; see the elliptic estimates in \cite{GilbargTrudinger2001}.
Consequently
\[
 \Delta_\delta(\check h_{ij}-\delta_{ij})
       =O_j(r^{-2-2q_0}),\qquad
 \Delta_\delta V_0=O_j(r^{-2-2q_0})
\]
for every $j$. Cut the fields off inside a fixed sphere. Their
Euclidean Laplacian sources are integrable and have a finite positive
moment of every order smaller than $2q_0-(n-2)$. Choose
\[
 0<\alpha<\min\{1,\,2q_0-(n-2)\}.
\]
Subtracting the Newton kernel monopole from the potential, the inner
annuli are bounded by the $\alpha$ moment and the outer annuli by
the displayed pointwise source bounds. Thus
\[
 \check h_{ij}-\delta_{ij}=A_{ij}r^{2-n}
                         +O_j(r^{2-n-\alpha}),\qquad
 V_0=-a_0r^{2-n}+O_j(r^{2-n-\alpha}).
\]
The difference from the Newton potential is an entire decaying
harmonic function and vanishes. Derivative estimates for the
remainders follow by the same scaled interior estimates.

Harmonic gauge now eliminates $A$. Its leading term is
\[
 0=\partial_j\check h_{ij}
       -\tfrac12\partial_i\check h_{jj}
       +O(|\check h-\delta|\,|\partial\check h|)
  =(2-n)r^{-n}
    \left(A_{ij}x^j-\tfrac12(\operatorname{tr}A)x_i\right)
    +o(r^{1-n}).
\]
Hence $A=(\operatorname{tr}A)I/2$, and taking the trace gives
$A=0$ because $n\ne2$. Finally the harmonic flux of $\lambda_0$
in $h_0$, with the outward domain normal $-\nu$ at $S$, gives
\begin{equation}\label{uniq:capacity-flux}
 (n-2)\omega a_0=\kappa_0\Area_{h_0}(S)>0.
\end{equation}
All changes between the Euclidean and metric fluxes have vanishing
limit by the decay just established.
\end{proof}

The next lemma specifies the inverse actually required at the point
of compactification. Only compact sources are used.

\begin{lemma}[Compact-source inverse and perturbation norm]
\label{uniq:compact-inverse}
Let $(M,\gamma)$ be connected, complete, without boundary, and have
one Euclidean coordinate end with compact complement. Suppose
$\gamma$ is uniformly positive, locally $W^{2,s}$, and smooth off
a compact set, where
\[
 \frac n2<s<n,\qquad p=\frac{ns}{n-s}>n.
\]
Suppose on its end $\gamma-\delta=O_j(r^{-q})$ for every $j$
and some $q>0$. Fix a compact set $K$ whose interior contains all
nonsmooth points, and a larger compact set $K_1$. For every
$0<b<n-2$ there is a unique operator $P_\gamma$ such that
\begin{equation}\label{uniq:inverse-estimate}
 \Delta_\gamma P_\gamma f=f,\qquad
 \|P_\gamma f\|_{W^{2,s}(K_1)}
  +\sup_{r\geq R} r^b
       \sum_{j=0}^2 r^j|D^jP_\gamma f|
       \leq C\|f\|_{L^s(K)}
\end{equation}
for $f\in L^s$ supported in $K$, with $P_\gamma f\to0$ at
infinity. The same bound holds for higher scaled derivatives on
the end, with a constant depending on the derivative order.

Write $\mathcal X=P_\gamma(L^s_K)$, endowed with
$\|v\|_{\mathcal X}=\|\Delta_\gamma v\|_{L^s(K)}$.
If $\gamma'$ agrees with $\gamma$ outside $K$ and is sufficiently
close in $W^{2,s}(K)$, then
\begin{equation}\label{uniq:operator-norm}
\begin{gathered}
 \left\| (\Delta_{\gamma'}-\Delta_\gamma)v
                 -a_n^{-1}R_{\gamma'}v\right\|_{L^s(K)}
 \leq C\left(\|\gamma'-\gamma\|_{W^{2,s}(K)}
                       +\|R_{\gamma'}\|_{L^s(K)}\right)
                     \|v\|_{\mathcal X},
 \\ a_n=\frac{4(n-1)}{n-2}.
\end{gathered}
\end{equation}
Here the curvature term is used when its support is contained in $K$.
\end{lemma}

\begin{proof}
Sobolev embedding gives continuous coefficients, first derivatives
in $L^p$, and uniform ellipticity on compact coordinate patches.
In nondivergence notation
$\Delta_\gamma=a^{ij}\partial_i\partial_j+B^i\partial_i$,
the principal coefficients are continuous and $B\in L^p$.
Freeze the principal coefficients on a small ball. The constant
coefficient $W^{2,s}$ estimate absorbs their small oscillation.
For the first-order term use
\[
 \|B\,Dv\|_{L^s}
 \leq \|B\|_{L^p}\|Dv\|_{L^n},\qquad
 \frac1s=\frac1p+\frac1n.
\]
Interpolation of $\|Dv\|_{L^n}$ between the local $W^{2,s}$
norm and a zeroth-order norm allows absorption, with an additional
zeroth-order constant. Cutoffs and a finite cover give the local
$W^{2,s}$ estimate. These estimates are uniform for coefficients
converging in $W^{2,s}$: their principal coefficients converge
uniformly, and the $L^p$ first-order coefficients are uniformly
equi-integrable. The same argument gives estimates for smooth
approximations and their Dirichlet problems.

For large $R$, direct differentiation gives
\begin{equation}\label{uniq:end-barrier}
 \Delta_\gamma r^{-b}
   =b(b+2-n)r^{-b-2}+O(r^{-b-2-q})<0
 \quad (r\geq R).
\end{equation}
Solve on expanding smooth truncations with zero outer Dirichlet
value, initially using smooth coefficient and source approximations.
The maximum principle on the exterior region bounds the solution
there by its supremum on $r=R$ times $C r^{-b}$; local estimates
also bound its scaled exterior derivatives. The compact-core
supremum is bounded by $C\|f\|_{L^s}$ uniformly in the
exhaustion and approximations. Otherwise divide by that diverging
supremum. The local estimates and the compact embedding of
$W^{2,s}$ into $C^0$ give a nonzero limit harmonic on all of $M$,
bounded by $Cr^{-b}$ at infinity. The divergence-form strong maximum
principle contradicts this limit. A maximizing point cannot escape
the fixed core because of the exterior barrier. Passing first in
the coefficient and source approximations on a truncation, and
then through exhaustion, proves existence and
\eqref{uniq:inverse-estimate}. The decaying homogeneous maximum
principle proves uniqueness, including for sign-changing solutions.
Thus $\mathcal X$ is a Banach space isometric to $L^s_K$.

For completeness, the perturbation estimate is an estimate between
these particular spaces, rather than an assertion about all weighted
sources. The second-order coefficient difference is bounded in
$C^0$ by $C\|\gamma'-\gamma\|_{W^{2,s}}$. The first-order
difference is bounded in $L^p$ by the same quantity, while
$Dv\in L^p(K)$ by \eqref{uniq:inverse-estimate}. Since
$p/2>s$, their product lies in $L^s(K)$. Finally
$R_{\gamma'}v\in L^s$ because $v\in C^0(K)$.
This proves \eqref{uniq:operator-norm}. All constants are for the
fixed sets and reference metric; one chooses the perturbation
size only after those constants have been fixed.
\end{proof}

\begin{lemma}[Conformal correction and its mass]
\label{uniq:correction-mass}
In Lemma~\ref{uniq:compact-inverse}, suppose in addition that
$R_\gamma=0$ weakly, that $\gamma$ has a finite ADM energy,
and that $q>(n-2)/2$. Choose $b$ so that
\begin{equation}\label{uniq:mass-exponents}
 \frac{n-2}{2}<b<n-2,\qquad b+q>n-2.
\end{equation}
Every sufficiently small compact $W^{2,s}$ perturbation $\gamma'$
has a unique small correction $v\in\mathcal X$ for which
\[
 U=1+v>0,\qquad
 \widehat\gamma=U^{4/(n-2)}\gamma',\qquad
 R_{\widehat\gamma}=0.
\]
The correction depends continuously on $\gamma'$ in $W^{2,s}$,
and differentiably along smooth one-parameter perturbations. Its
energy satisfies
\begin{equation}\label{uniq:mass-formula}
 \mathcal E(\widehat\gamma)
 =\mathcal E(\gamma)
  -\frac{1}{2(n-1)\omega}
      \int_M R_{\gamma'}(1+v)\,\dd V_{\gamma'}.
\end{equation}
In particular the energies converge under compact smoothing in
$W^{2,s}$ followed by this correction.
\end{lemma}

\begin{proof}
The conformal scalar-curvature equation is
\[
 -a_n\Delta_{\gamma'}v+R_{\gamma'}(1+v)=0.
\]
Equivalently, on $\mathcal X$,
\begin{equation}\label{uniq:neumann-correction}
 v=P_\gamma\left[
 a_n^{-1}R_{\gamma'}
 -(\Delta_{\gamma'}-\Delta_\gamma)v
 +a_n^{-1}R_{\gamma'}v\right].
\end{equation}
Curvature is a continuous map from positive $W^{2,s}$ metrics to
$L^s$: its second-order terms have continuous coefficients, and
its quadratic first-order terms are in $L^{p/2}\subset L^s$.
Thus $R_{\gamma'}$ tends to zero in $L^s$. By
\eqref{uniq:operator-norm}, the operator on the right involving
$v$ has norm less than $1/2$ after shrinking the metric
neighborhood. Its Neumann inverse gives the solution and its stated
dependence. Estimate~\eqref{uniq:inverse-estimate} makes
$\|v\|_{C^0}<1/2$, so $U$ is positive.

Outside $K$, $v$ is $\gamma$-harmonic. The divergence theorem,
valid also for the weak equation on the compact part, therefore
gives the finite flux
\begin{equation}\label{uniq:correction-flux}
 F(v):=\lim_{r\to\infty}\int_{S_r}\partial_{\nu_\gamma}v
                 \,\dd A_\gamma
       =a_n^{-1}\int_M R_{\gamma'}(1+v)\,\dd V_{\gamma'}.
\end{equation}
The usual conformal ADM calculation in the fixed end chart gives
\[
 \mathcal E(U^{4/(n-2)}\gamma')-\mathcal E(\gamma')
       =-\frac{2}{(n-2)\omega}F(v).
\]
Here the errors from $vDv$ have integral
$O(r^{n-2-2b})$, and those from the background metric and its first
derivative have integral $O(r^{n-2-q-b})$. Both vanish by
\eqref{uniq:mass-exponents}. Also $\gamma'=\gamma$ at infinity,
so its uncorrected ADM energy is unchanged. This proves
\eqref{uniq:mass-formula}, including existence of the corrected
ADM limit. Continuity and differentiation follow from the compact
integral, the operator convergence in
\eqref{uniq:neumann-correction}, and $W^{2,s}\hookrightarrow C^0$.
\end{proof}

\begin{lemma}[Smooth nonnegative mass from the numerical theorem]
\label{uniq:smooth-nonnegative-mass}
Let $(M,\gamma)$ be a smooth connected oriented complete manifold
without boundary, with one Euclidean coordinate end and compact
complement, scalar curvature zero, and finite ADM energy. Assume that $\gamma-\delta=O_j(r^{-q})$ for every integer $j\geq0$,
with $q>(n-2)/2$. Then $\mathcal E(\gamma)\geq0$.
\end{lemma}

\begin{proof}
Fix a smooth point $o$. The local Dirichlet Green function for
$-\Delta_\gamma$ in a small ball has a positive fundamental pole
at $o$. Cut it off inside that ball. Its distributional Laplacian
is $-\delta_o$ plus a smooth compact source supported in an annulus
away from $o$. Subtract the decaying solution of that compact
source given by Lemma~\ref{uniq:compact-inverse}. The result $G$
is harmonic off $o$, tends to zero at infinity, and has the same
positive pole. The maximum principle on the region between a
small sphere about $o$ and a large outer sphere gives $G\geq0$;
the strong principle gives $G>0$. In normal distance $d$ from $o$,
the local fundamental solution estimates give, with $c>0$,
\[
 G=c\,d^{2-n}(1+o(1)),\qquad
 \partial_dG=(2-n)c\,d^{1-n}(1+o(1)),
\]
uniformly in angle. These local estimates follow by rescaling the
smooth elliptic coefficients in the Green equation. At the retained
end $G=O_2(r^{-b})$ for any $0<b<n-2$, by
Lemma~\ref{uniq:compact-inverse} and exterior estimates; its
harmonic flux there is finite.

For fixed $\epsilon>0$, set
\[
 \gamma_\epsilon=(1+\epsilon G)^{4/(n-2)}\gamma
 \quad\hbox{on }M\setminus\{o\}.
\]
This metric is scalar flat. On a sufficiently small sphere $d=t$,
the normal toward the old end is the increasing-$d$ normal, and
its conformal mean curvature has the sign of
\[
 H_\gamma+\frac{2(n-1)}{n-2}
        \partial_d\log(1+\epsilon G)
 =\frac{n-1}{t}-\frac{2(n-1)}{t}+o(t^{-1})<0.
\]
For this fixed $\epsilon$, choose such a sphere and discard its
inside. The remaining manifold is a connected smooth exterior,
complete including its compact boundary, with precisely the old
coordinate end. Give it second tensor and electric field zero.
It has zero matter constraints, a strictly future trapped boundary,
and the required $O_2$ metric decay with any fixed exponent
between $(n-2)/2$ and $\min\{q,n-2\}$. Its energy exists:
choose $b$ additionally so that $b+q>n-2$, and use the finite
harmonic flux and the ADM calculation of
Lemma~\ref{uniq:correction-mass}.

The numerical Theorem~\ref{thm:numerical}, at $Q=0$, applies to
this exterior. In particular it proves that its energy is positive;
timelikeness and a positive full-cut infimum are conclusions of
that theorem, not additional assumptions here. The conformal mass
formula gives
\[
 \mathcal E(\gamma_\epsilon)
 =\mathcal E(\gamma)
   -\frac{2\epsilon}{(n-2)\omega}
       \lim_{r\to\infty}\int_{S_r}\partial_{\nu_\gamma}G
                              \,\dd A_\gamma
 \longrightarrow\mathcal E(\gamma).
\]
Send $\epsilon\downarrow0$, choosing the small inner sphere only
after each $\epsilon$ is fixed. This proves the assertion. The
argument uses only the already established numerical theorem in
dimension $n$ and imposes no spin assumption.
\end{proof}

\begin{lemma}[Rigidity of the compactified zero-mass metric]
\label{uniq:rough-zero-mass}
Let $(M,\gamma)$ be oriented and satisfy the hypotheses of
Lemma~\ref{uniq:compact-inverse}, with $R_\gamma=0$ weakly,
$\mathcal E(\gamma)=0$, and $q>(n-2)/2$. Suppose $\gamma$ is
smooth except at one point and possibly along a smooth compact
hypersurface where it is $C^{2,\beta}$. Then $(M,\gamma)$ is
Euclidean space after passage to smooth harmonic coordinates.
\end{lemma}

\begin{proof}
Fix $K$ containing all nonsmooth points. A finite coordinate
partition, mollification, and a cutoff within $K$ give smooth
positive metrics $\gamma_j$ agreeing with $\gamma$ off $K$
and converging in $W^{2,s}$ and uniformly. Positivity follows
from uniform convergence and the compact lower ellipticity bound.
By Lemma~\ref{uniq:correction-mass}, they have positive conformal
corrections $U_j\to1$. The corrected metrics are smooth by elliptic
regularity, complete because the factors are uniformly bounded
above and below, scalar flat, and have one retained coordinate
end with decay exponent $\min\{q,b\}>(n-2)/2$. Their energies converge to
zero by \eqref{uniq:mass-formula} and are nonnegative by
Lemma~\ref{uniq:smooth-nonnegative-mass}.

More generally, let $H$ be an arbitrary smooth symmetric tensor
compactly supported in a smooth patch of $\gamma$, and set
$\gamma_t=\gamma+tH$ for $|t|$ sufficiently small. Enlarge $K$
once to contain $\supp H$, and choose its reference inverse before
varying $t$. Use this same $P_\gamma$, this same $K$, and
\eqref{uniq:neumann-correction} to obtain $U_t=1+v_t$ and
scalar-flat $\widehat\gamma_t$. For each fixed $t$, approximate
$\gamma_t$ smoothly in $W^{2,s}$, keeping its exterior fixed,
and correct using that same reference inverse. The corrected
smooth metrics converge to $\widehat\gamma_t$, and their
energies converge by \eqref{uniq:mass-formula}. Consequently
\[
 \mathcal E(\widehat\gamma_t)\geq0
 \quad\hbox{for both signs of }t,
 \qquad \mathcal E(\widehat\gamma_0)=0.
\]
Only after taking the smoothing limit at each fixed $t$ do we
differentiate this two-sided inequality at $t=0$.

Since $R_\gamma=0$ and $U_0=1$, differentiation of
\eqref{uniq:mass-formula} gives
\begin{equation}\label{uniq:mass-derivative}
 0=\left.\frac{\dd}{\dd t}\right|_0
          \mathcal E(\widehat\gamma_t)
  =-\frac{1}{2(n-1)\omega}\int_M R'_\gamma[H]\,\dd V_\gamma
  =\frac{1}{2(n-1)\omega}
          \int_M\langle\Ric_\gamma,H\rangle\,\dd V_\gamma.
\end{equation}
The last identity follows by integrating
$R'[H]=\Div\Div H-\Delta\tr H-\langle\Ric,H\rangle$;
the support lies in a smooth interior patch, so no boundary terms
occur. Thus $\Ric_\gamma=0$ on every smooth patch. It vanishes
across the hypersurface by its $C^{2,\beta}$ regularity as well.
At the exceptional point, $D^2\gamma\in L^s$ and
$D\gamma\,D\gamma\in L^{p/2}\subset L^s$ imply that
distributional Ricci curvature has $L^s$ coefficients. There is
therefore no point-supported curvature term, and $\Ric_\gamma=0$
weakly throughout $M$.

We detail the regularity upgrade, so that comparison geometry is
applied to a smooth metric. Fix $n/2<s<n$ and
$p=ns/(n-s)>n$. In an original chart the scalar Laplacian has
the form $a^{ij}\partial_i\partial_j+B^i\partial_i$,
where $a$ is continuous, $Da,B\in L^p$, and $DB\in L^s$.
Indeed $DB$ consists of bounded coefficients times $D^2\gamma$
and quadratic $D\gamma$ terms in $L^{p/2}\subset L^s$.
Normalize $a(x_0)=I$ by a constant linear change and rescale
$x=x_0+r z$. The rescaled drift $B_r(z)=rB(x_0+r z)$ satisfies
\[
 \|a(x_0+r\,\cdot)-I\|_{\infty}\longrightarrow0,
 \qquad
 \|B_r\|_{L^p(B_2)}
   =r^{1-n/p}\|B\|_{L^p(B_{2r}(x_0))}\longrightarrow0.
\]
Let $G_D$ be the flat Dirichlet inverse on $B_1$. On
$\mathcal D_p=W^{2,p}(B_1)\cap W^{1,p}_0(B_1)$ the equations
for harmonic replacements $y^i=z^i+v^i$ are
\[
 v^i=G_D\bigl[-B_r^i+(I-a_r):D^2v^i-B_r\cdot Dv^i\bigr].
\]
The operator involving $v^i$ has norm at most
$C_p(\|I-a_r\|_\infty+\|B_r\|_p)$, because
$\|Dv\|_\infty\leq C_p\|v\|_{W^{2,p}}$. Choose $r$ so this
norm is less than $1/2$. The Neumann series gives existence,
uniqueness in $\mathcal D_p$, and
$\|v^i\|_{W^{2,p}}\leq C_p\|B_r^i\|_p$. A further decrease
of $r$ makes $\|Dy-I\|_\infty<1/2$. Integrating along straight
segments in $B_1$ makes $y$ injective and bilipschitz, so it gives
harmonic coordinates on a smaller ball.

Approximate the metric smoothly in $W^{2,s}$ on the rescaled
ball. Then $a_j\to a$ uniformly, $B_j\to B$ in $L^p$, and
$DB_j\to DB$ in $L^s$. The same contraction gives harmonic
coordinates $y_j\to y$ strongly in $W^{2,p}$ and $C^1$.
For $w_j=\partial_\ell y_j$, differentiation gives
\[
 a_j:D^2w_j
 =-(\partial_\ell a_j):D^2y_j
   -(\partial_\ell B_j)\cdot Dy_j-B_j\cdot Dw_j.
\]
The right side is bounded in $L^s$: its first and last products
are in $L^{p/2}\subset L^s$, and its middle product is in
$L^s$. Local frozen-coefficient estimates give uniform
$W^{2,s}$ bounds for $w_j$ on smaller balls. Applied to differences,
these estimates and the already strong $W^{2,p}$ convergence give
strong local $W^{3,s}$ convergence of $y_j$ to $y$.

The inverse chart $x=y^{-1}$ is $C^1$ bilipschitz and
$W^{2,p}\cap W^{3,s}$. This follows by twice and three times
differentiating $y(x)=\id$: the second derivatives of $x$ are
bounded factors times $D^2y\circ x$, and its third derivatives
are bounded factors times $D^3y\circ x$ or
$(D^2y\circ x)^2$. The latter belongs to $L^{p/2}\subset L^s$.
The transformed metric remains $W^{2,s}$ for the same reason:
its second derivatives contain $D^2\gamma$ and $D^3x$ in $L^s$,
and products of two factors among $D\gamma$ and $D^2x$ in
$L^{p/2}$. On a fixed smaller image chart the smooth transformed
metrics converge strongly in $W^{2,s}$. Composition convergence
follows by approximation of the composed functions by smooth
ones and the uniform Jacobian bounds; the other factors converge
in the displayed norms. Since Ricci curvature is continuous
from uniformly elliptic $W^{2,s}$ metrics to $L^s$, the smooth
Ricci transformation laws pass to this limit. Thus the weakly
Ricci-flat metric satisfies, in these harmonic coordinates,
\[
 a^{ab}\partial_a\partial_b\gamma_{ij}
        =Q_{ij}(\gamma^{-1},D\gamma),
\]
where $Q$ is quadratic in $D\gamma$ and belongs to $L^t$ with
$t=p/2>s$.

To obtain actual higher regularity from this equation, take a
cutoff $\eta$ supported in a small ball and put
$u=\eta\gamma_{ij}$. It has zero trace and satisfies
\[
 a:D^2u=\eta Q_{ij}
       +2a(D\eta,D\gamma_{ij})+(a:D^2\eta)\gamma_{ij}
       =:F\in L^t,
\]
using $D\gamma\in L^p\subset L^t$ and boundedness of
$\gamma$. Choose the ball so that, after normalizing $a$ at
its center, the oscillation is small enough for the Dirichlet
Neumann inverse of $a:D^2$ at both exponents $s$ and $t$.
The inverse at $t$ gives a zero-trace solution
$v\in W^{2,t}$ of $a:D^2v=F$. Since $t>s$, both $u$ and
$v$ belong to $W^{2,s}$; their difference has zero trace and
solves the homogeneous equation. Uniqueness of the same inverse
at $s$ gives $u=v$. Thus the metric actually gains $W^{2,t}$
regularity on the smaller ball, rather than merely satisfying a
conditional higher-norm estimate.

As long as $s<n$, the improved exponent is
$t=ns/[2(n-s)]$, so $1/t=2/s-2/n$. Since $s>n/2$, finitely
many iterations reach an exponent greater than $n$, taking
slightly smaller exponents if an endpoint is met. The metric
is then $C^{1,\beta}$, its quadratic source is $C^{0,\beta}$,
and Schauder gives $C^{2,\beta}$ followed by smoothness.
The classical harmonic-coordinate regularity theorem
\cite{DeTurckKazdan1981} is used only for this final smooth
continuation, after the initial Sobolev steps just proved.

The resulting smooth metric is complete and Ricci flat. Its single
Euclidean end gives asymptotic volume ratio one: the compact part
has finite volume and the metric error tends to zero, so distance
and volume comparison with exterior Euclidean annuli gives this
limit for balls about any fixed point. Bishop--Gromov comparison
and its equality case \cite{Wei2007} now apply. The volume ratio
is nonincreasing, has limit one both at radius zero and at infinity,
and hence is identically one. The equality case identifies the
complete manifold with $\R^n$.
\end{proof}

\begin{proof}[Proof of Proposition~\ref{uniq:vacuum-uniqueness}]
At $S$, \eqref{uniq:static-equations} gives
$\Hess\lambda_0=0$. Since $\partial_\nu\lambda_0=\kappa_0>0$,
the tangential Hessian identity shows that $S$ is totally geodesic.
In Gaussian coordinates, reflect $h_0$ evenly and $\lambda_0$
oddly. The vanishing first normal derivative of the tangential
metric and the vanishing second normal derivative of $\lambda_0$
make these extensions $C^{2,\beta}$ for $0<\beta<1$. The static
equations hold across the seam by continuity. On the resulting
two-ended double set
\begin{equation}\label{uniq:double-metric}
 \gamma=\left(\frac{1+\lambda_0}{2}\right)^{4/(n-2)}h_0,
\end{equation}
where $\lambda_0$ is now the odd extension. The factor is positive
at every finite point and on the seam. The conformal scalar law,
$R_{h_0}=0$, and $\Delta_{h_0}\lambda_0=0$ show that
$R_\gamma=0$, also across the seam.

On the retained end,
\[
 \gamma=
 \left(\frac{(1+\lambda_0)^2}{4\lambda_0}\right)^{2/(n-2)}
      \check h,
 \qquad
 \frac{(1+\lambda_0)^2}{4\lambda_0}
     =1+\frac{(1-\lambda_0)^2}{4\lambda_0}.
\]
Lemma~\ref{uniq:end-expansion} therefore implies a decay strictly
faster than $r^{2-n}$, with all derivatives, and zero ADM energy.
On the other end, expressed using the positive lapse on its original
copy,
\[
 \gamma=
 \left(\frac{(1-\lambda_0)^2}{4\lambda_0}\right)^{2/(n-2)}
       \check h
 =\left(\frac{a_0}{2}\right)^{4/(n-2)}r^{-4}
       \bigl(\delta+O_j(r^{-\alpha'})\bigr)
\]
for some $\alpha'>0$, decreasing it if necessary. In inverted
coordinates $y=x/|x|^2$ this is a positive constant Euclidean
metric plus an error satisfying
\begin{equation}\label{uniq:point-estimate}
 |D^j(\gamma-c\delta)|\leq C_j|y|^{\alpha'-j},
 \qquad c=\left(\frac{a_0}{2}\right)^{4/(n-2)}>0.
\end{equation}
Fill $y=0$ with that value. Take $0<\alpha'<2$ and then fix
\begin{equation}\label{uniq:point-exponent}
 \frac n2<s<\min\left\{n,\frac{n}{2-\alpha'}\right\}.
\end{equation}
This is a nonempty interval. Estimate~\eqref{uniq:point-estimate}
gives $W^{2,s}$ coefficients, since
$s(\alpha'-2)+n>0$. They are continuous and positive at the
point. There is no distributional derivative concentrated there:
the boundary terms obtained by integrating the first derivatives
over $|y|=\epsilon$ tend to zero by
\eqref{uniq:point-estimate}. Equivalently these are genuine
$W^{2,s}$ extensions. Scalar curvature is thus $L^s$, and its
vanishing off the point implies weak vanishing everywhere.

The filled double is connected, oriented, complete, without
boundary, and has just the retained coordinate end with compact
complement. Completeness near the filled point follows from its
uniform ellipticity; it follows at the seam from its regular
positive metric. It satisfies Lemma~\ref{uniq:rough-zero-mass},
so it is Euclidean space.

The seam is a connected compact embedded totally umbilic
hypersurface in this Euclidean space. Indeed its original second
fundamental form is zero, and the conformal change
\eqref{uniq:double-metric} adds a constant multiple of its induced
metric, with nonzero coefficient because $\kappa_0>0$. A connected
Euclidean totally umbilic hypersurface with nonzero principal
curvature is contained in a sphere: the Codazzi equation makes
that curvature constant, and differentiation of its center
$x-\nu/\kappa_{\mathrm{Euc}}$ gives zero. Compactness makes it
the whole sphere. The retained component is its unbounded side.
After a translation denote the sphere radius by $R$.

On that component $h_0=U^{4/(n-2)}\delta$, with
$U=2/(1+\lambda_0)$. Scalar flatness makes $U$ Euclidean harmonic,
and it has boundary value two and limit one at infinity. The
exterior Dirichlet maximum principle identifies it uniquely as
$1+R^{n-2}\rho^{2-n}$. This proves
\eqref{uniq:schwarzschild-pair}. The smoothness of the identification
at the seam follows either from its harmonic coordinates or from
normal geodesic collars of the two smooth metrics.
\end{proof}
